\chapter{Implementación computacional}
\label{app:codigo}
\singlespacing
\lstset{basicstyle=\ttfamily\scriptsize,breaklines=true,numbersep=4pt,frame=none}

Se presentan los programas empleados en el Capítulo~\ref{chap:resultados}. Los listados A.1--A.12 documentan los experimentos iniciales del proyecto; el A.8 presenta ahora el código corregido para nuevas ejecuciones, mientras que el archivo histórico original se conserva aparte para la trazabilidad de los CSV publicados. Los experimentos adicionales y las correcciones del protocolo se reproducen en el cuaderno incluido en \texttt{cuadernos/}; las secciones añadidas al final de este apéndice detallan su procedencia y los archivos necesarios. Los métodos comparten la ecuación de Poisson cuando corresponde, pero las nuevas celdas deben ejecutarse después de las definiciones indicadas en el cuaderno.

El cuaderno computacional se encuentra en el
\href{https://github.com/MichaelMobius/tesispucp}{repositorio público \texttt{tesispucp} de GitHub}.
Los resultados originales y los nuevos registros de ejecución se conservan por separado. El barrido 2D adicional utilizó el perfil \texttt{experimento}, mientras que la reproducción principal de Shin empleó el perfil \texttt{articulo}. El código diferencia los errores relativos
computados en los ejemplos propios de los errores RMS discretos no relativos
utilizados para contrastar las figuras del artículo de
Shin et al.~\cite{shin2020convergence}.

\noindent\textbf{Dependencias.} Python 3, PyTorch~\cite{paszke2019pytorch}, NumPy, SciPy, Matplotlib y pandas;
para la comparación con métodos clásicos, también \texttt{scikit-fem}~\cite{gustafsson2020scikitfem} y
\texttt{FiPy}~\cite{guyer2009fipy}. La precisión doble se usa en las reproducciones de Shin y
en la comparación de regularizaciones; los ejemplos iniciales emplean,
según su configuración, precisión simple.

\section{Regresión simple}
\label{app:experimento1}
\lstinputlisting[language=Python,caption={Regresión con red neuronal simple},label={lst:anexo1}]{codigos/A01_regresion.py}

\section{Ecuación logística}
\label{app:experimento2}
\lstinputlisting[language=Python,caption={PINN para la ecuación logística},label={lst:anexo2}]{codigos/A02_logistica.py}

\section{Poisson 1D con observaciones: solución seno}
\label{app:experimento3}
\lstinputlisting[language=Python,caption={Poisson 1D con datos: solución seno},label={lst:anexo3}]{codigos/A03_poisson1d_seno_datos.py}

\section{Poisson 1D con observaciones: tangente hiperbólica}
\label{app:experimento4}
\lstinputlisting[language=Python,caption={Poisson 1D con datos: tangente hiperbólica},label={lst:anexo4}]{codigos/A04_poisson1d_tanh_datos.py}

\section{Poisson 1D sin observaciones interiores}
\label{app:experimento5}
\lstinputlisting[language=Python,caption={Poisson 1D sin observaciones},label={lst:anexo5}]{codigos/A05_poisson1d_tanh_sin_datos.py}

\section{Comparación de regularizadores}
\label{app:experimento6}
\lstinputlisting[language=Python,caption={PINN estándar, penalización media y LIPR},label={lst:anexo6}]{codigos/A06_regularizaciones.py}

\section{Reproducción de Shin: Dirichlet débil}
\label{app:experimento7}
\lstinputlisting[language=Python,caption={Poisson 1D: reproducción de la Figura 3 de Shin},label={lst:anexo7}]{codigos/A07_shin_tanh.py}

\section{Reproducción de Shin: Dirichlet fuerte}
\label{app:experimento8}
\noindent El Listado~\ref{lst:anexo8} es una implementación corregida, destinada a \emph{ejecuciones futuras}: fija el peso LIPR mediante el total $N_f$ y escribe en un directorio separado. Los CSV históricos de la Tabla~\ref{tab:shin_diagnostico_nuevo} y de la reproducción principal \emph{no} se recalcularon con este archivo. La versión original permanece como \texttt{codigos/A08\_shin\_oscilatorio\_historico.py}, pero no se repite integralmente para evitar duplicación.
\lstinputlisting[language=Python,caption={Protocolo oscilatorio corregido para nuevas ejecuciones},label={lst:anexo8}]{codigos/A08_shin_oscilatorio.py}

\section{Regularidad Hölder y cobertura}
\label{app:experimento9}
Se estima la seminorma de H\"older a partir de una malla auxiliar finita. El máximo discreto no certifica el supremo continuo de los teoremas.

\lstinputlisting[language=Python,caption={Aproximación discreta de seminormas Hölder},label={lst:anexo9}]{codigos/A09_holder.py}

\section{Poisson 2D con observaciones}
\label{app:experimento10}
\lstinputlisting[language=Python,caption={PINN Poisson 2D con observaciones},label={lst:anexo10}]{codigos/A10_poisson2d_datos.py}

\section{Poisson 2D sin observaciones}
\label{app:experimento11}
\lstinputlisting[language=Python,caption={PINN Poisson 2D sin observaciones},label={lst:anexo11}]{codigos/A11_poisson2d_sin_datos.py}

\section{Comparación FEM, FVM y PINN}
\label{app:experimento12}
Los errores se evalúan sobre una cuadrícula común.

\lstinputlisting[language=Python,caption={Comparación numérica FEM, FVM y PINN},label={lst:anexo12}]{codigos/A12_fem_fvm_pinn.py}



\section{Datos, entorno y trazabilidad del protocolo revisado}
\label{app:protocolo_revisado}

\paragraph{Cuaderno y fuentes de datos.}
El cuaderno de Colab se adjunta al proyecto como archivo ejecutable.
Los datos numéricos utilizados para actualizar los Capítulos IV y V se adjuntan en la carpeta \texttt{resultados/}, sin sustituir los CSV históricos. Los archivos clave son:
\begin{itemize}
\item \texttt{convergencia\_poisson2d\_revisada.csv}: una fila por cada una de las $24$ ejecuciones, con número de puntos, semilla, tipo de contorno, errores, residuos, defectos de frontera, tiempos y firma de configuración;
\item \texttt{resumen\_convergencia\_poisson2d.csv}: medianas y desviaciones estándar de las tres semillas;
\item \texttt{estabilidad\_modelos\_originales\_2d.csv}: reevaluación, en una malla distinta, de los dos experimentos 2D originales;
\item \texttt{diagnostico\_oscilatorio\_controlado.csv}: comparación controlada de lote completo y minilote en el problema 1D de alta frecuencia;
\item \texttt{resultados\_shin\_oscilatorio.csv}: tabla principal de la reproducción oscilatoria con la versión de función objetivo registrada;
\item \texttt{comparacion\_fem\_fvm\_pinn.csv}: resultados del benchmark con $N=30$ y tiempos de ejecución;
\item \texttt{entorno\_revision.json}: versiones del entorno histórico; el archivo \texttt{HISTORICO\_manifiesto\_sha256.csv} es propio de otra ejecución y no certifica este depósito; \texttt{MANIFIESTO\_DEPOSITO\_SHA256.csv} es el inventario generado para los archivos que se entregan ahora.
\end{itemize}

El entorno registrado fue Python 3.13.16, PyTorch 2.11.0 con CUDA 13.0, NumPy 2.1.3, SciPy 1.16.3, FiPy 4.0.3 y scikit-fem 12.0.2, sobre una GPU NVIDIA A100-SXM4-40GB. En el barrido 2D se empleó \texttt{torch.float32} y en Shin \texttt{torch.float64}. Se adjuntan los $24$ modelos de barrido 2D en \texttt{modelos\_2d\_barrido/}, los ocho modelos de diagnóstico oscilatorio en \texttt{diagnostico\_shin\_modelos/} y los checkpoints disponibles en \texttt{checkpoints\_shin/}. No debe asumirse que un archivo de pesos histórico es intercambiable entre versiones de código con distinta función objetivo.

\paragraph{Discretización y precisión numérica.}
En las reproducciones unidimensionales de Shin, la función \texttt{linspace} genera $N_f$ nodos desde $-1$ hasta $1$ e incluye los extremos para $N_f\geq2$, y el caso $N_f=1$ utiliza el origen; por ello no todos esos puntos pertenecen al dominio abierto. En Poisson 2D, $N=75$ representa subdivisiones por eje: $(N-1)^2=5\,476$ puntos interiores y $4(N+1)=304$ entradas de frontera, con cuatro esquinas duplicadas. El peso de penalización de frontera de los experimentos 2D originales es $10$. La ejecución 2D usa derivadas automáticas segundas en precisión simple (\texttt{float32}); el protocolo unidimensional de Shin usa \texttt{float64}. Los efectos de redondeo en segunda derivada no se cuantificaron mediante un experimento de sensibilidad de precisión.

\paragraph{Definición de época y objetivo LIPR.}
En la revisión se fijó el peso $\lambda_R=N_f^{-3/2}$ a partir del tamaño total de los puntos de colocación, incluso cuando se entrena con un lote más pequeño. Cuando se utilizan minilotes, la cantidad \texttt{max} del regularizador se evalúa en el lote y no es idéntica al máximo en todo el conjunto. En este documento, \emph{actualización Adam} designa una aplicación del optimizador a un lote; no debe confundirse con una pasada completa por los $N_f$ ejemplos. El protocolo específico y sus firmas se encuentran en las celdas que generan los CSV y los checkpoints.

\paragraph{Rectas auxiliares y regularizadores discretos.}
En las figuras de Shin, la recta de pendiente $-1$ se coloca verticalmente con la mediana de los errores positivos del barrido; es una guía visual dependiente de los datos, no un ajuste estadístico de pendiente ni una tasa demostrada. Para $\alpha=1$, el máximo de la derivada del operador y la seminorma Lipschitz son equivalentes como magnitudes continuas bajo regularidad suficiente, pero los máximos discretos LIPR y H\"older se calculan en muestras diferentes. El caso $\alpha=1/2$ se evaluó, sin un entrenamiento separado.

\paragraph{Métricas en Poisson 2D.}
Para $Q$ malla de prueba y $Q_\partial$ puntos de su frontera, los valores reportados son
\begin{align*}
 J_{\mathrm{PDE,test}}&=\frac{1}{|Q|}\sum_{z\in Q}|\Delta h_\theta(z)-f(z)|^2,\\
 R_{\infty,Q}&=\max_{z\in Q}|\Delta h_\theta(z)-f(z)|,\\
 E_{\infty,Q}&=\max_{z\in Q}|h_\theta(z)-u^*(z)|,\\
 D_{\partial,Q}&=\max_{z\in Q_\partial}|h_\theta(z)-g(z)|.
\end{align*}
La columna \texttt{cota\_estabilidad\_evaluada} almacena $B_Q=R_{\infty,Q}/2+D_{\partial,Q}$. No es una cota certificada: el máximo de un residuo sobre una malla es, en general, \emph{menor o igual} que su supremo continuo, y tampoco certifica el supremo del error. Los errores $L^2$ y $H^1$ informados se estiman discretamente y se normalizan respecto de la solución analítica.

\paragraph{Extracción de implementaciones.}
Para facilitar la auditoría sin duplicar todo el cuaderno, los siguientes listados corresponden a sus celdas de revisión. Se ejecutan dentro del cuaderno después de sus definiciones anteriores, y no constituyen programas independientes.
\lstinputlisting[language=Python,caption={Dirichlet fuerte y métricas discretas de estabilidad: extracto del cuaderno},label={lst:anexo13}]{codigos/A13_metricas_estabilidad_2d.py}
\lstinputlisting[language=Python,caption={Entrenamiento revisado y registro reproducible del barrido 2D: extracto del cuaderno},label={lst:anexo14}]{codigos/A14_barrido_dirichlet_2d.py}
\lstinputlisting[language=Python,caption={Diagnóstico controlado del caso oscilatorio: extracto del cuaderno},label={lst:anexo15}]{codigos/A15_diagnostico_oscilatorio.py}

\paragraph{Reproducibilidad y limitaciones.}
Los valores presentados se refieren a la ejecución registrada, no a una certificación matemática. Los cambios de precisión numérica, biblioteca, plataforma, semilla o función objetivo pueden modificar los resultados. La combinación de CSV, gráficos, parámetros, checkpoints y cuaderno permite inspeccionar las correspondencias entre datos y figuras; para verificar un número específico se debe ejecutar la celda correspondiente con su configuración completa y comparar la huella del archivo exportado.


\section{Protocolo de comprobaci\'on de los comparadores aproximados (A--C)}
\label{app:tanh_abc}

El cuaderno ejecutable y la nota de procedencia se adjuntan al
directorio de reproducibilidad. Las medianas reportadas est\'an
en un CSV de la carpeta de resultados; su nombre y las reglas
de interpretaci\'on, figura en la nota de procedencia. Las figuras
del estudio se reconstruyeron a partir de las medianas reportadas,
no de registros individuales inventados.

El depósito incluye el CSV individual recuperado del Colab (16 filas),
el resumen original, el entorno y las cuatro fuentes de checkpoints del
protocolo A. Los archivos originales de Drive permiten contrastar los
resultados por semilla; los modelos finales y los historiales por etapa
no se han incorporado íntegramente a esta copia. Para una auditoría de
los entrenamientos finales debe consultarse el archivo original en
Google Drive, identificado en la nota de procedencia. El manifiesto
SHA-256 del depósito se genera sobre los archivos efectivamente
entregados, y el manifiesto antiguo se conserva como histórico.
